\documentclass{article}

\usepackage[series={A},noend,nofamiliar,noeledsec]{reledmac}

\firstlinenum{1}
\linenumincrement{1}
\numberstanzatrue
\Xstanza
\Xstanzaseparator{.}


\Xnumberonlyfirstinline
\setstanzaindents{0,0,0,0,0,0,0,0,0}


\begin{document}

\beginnumbering

\stanza
\edlabel{I1}\edtext{Vorrei}{\xxref{I1}{I2}\lemma{Vorrei \dots umani}\Afootnote{Reference inside one stanza}}
cantar quel memorando sdegno&
\edtext{ch’infiammò}{\Afootnote{Reference inside the same stanza, with no crossref}} già ne’ fieri petti umani\edlabel{I2}&
un’infelice e vil secchia di legno&
che tolsero a i Petroni i Gemignani.&
\edlabel{I3}\edtext{Febo}{\xxref{I3}{I4}\lemma{Febo \dots biondo}\Afootnote{Reference crossing two stanzas}},
che mi raggiri entro lo ’ngegno&
l’orribil guerra e gl’accidenti strani,&
tu che sai poetar servimi d’aio&
\edlabel{I5}\edtext{e tiemmi}{\lemma{e tiemmi \dots Elena}\xxref{I5}{I6}\Afootnote{Another lemma cross stanza}}
per le maniche del saio.\&

\stanza
E tu nipote del rettor del mondo&
del generoso Carlo ultimo figlio,&
ch’in giovinetta guancia e ’n capel biondo\edlabel{I4}&
copri canuto senno, alto consiglio,&
se da gli studi tuoi di maggior pondo&
volgi talor per ricrearti il ciglio,&
vedrai, s’al cantar mio porgi l’orecchia,&
\edlabel{I6}Elena trasformarsi in una secchia.\&

\stanza
Line 1 &
Line 2 &
Line 3 &
Line 4 &
Line 5 &
Line 6 &
Line 7 &
\edtext{Line 8}{\Afootnote{A note on line 8, but not the same stanza}}\&

\endnumbering

\noindent I3 stanza : \xstanzaref{I3}\\
I4 stanza : \xstanzaref{I4}
\newpage

\Xpstart
\Xstanza[][false]
\numberpstarttrue
\lineation{pstart}
\beginnumbering

\pstart
\edlabel{Ib1}\edtext{Vorrei}{\xxref{Ib1}{Ib2}\lemma{Vorrei \dots umani}\Afootnote{Reference inside one pstart}}
cantar quel memorando sdegno\\
\edtext{ch’infiammò}{\Afootnote{Reference inside the same pstart, with no crossref}} già ne’ fieri petti umani\edlabel{Ib2}\\
un’infelice e vil secchia di legno\\
che tolsero a i Petroni i Gemignani.\\
\edlabel{Ib3}\edtext{Febo}{\xxref{Ib3}{Ib4}\lemma{Febo \dots biondo}\Afootnote{Reference crossing two pstart}},
che mi raggiri entro lo ’ngegno\\
l’orribil guerra e gl’accidenti strani,\\
tu che sai poetar servimi d’aio\\\
\edlabel{Ib5}\edtext{e tiemmi}{\lemma{e tiemmi \dots Elena}\xxref{Ib5}{Ib6}\Afootnote{Another lemma cross pstart}}  per le maniche del saio.
\pend
\pstart
E tu nipote del rettor del mondo\\
del generoso Carlo ultimo figlio,\\
ch’in giovinetta guancia e ’n capel biondo\edlabel{Ib4}\\
copri canuto senno, alto consiglio,\\
se da gli studi tuoi di maggior pondo\\
volgi talor per ricrearti il ciglio,\\
vedrai, s’al cantar mio porgi l’orecchia,\\
\edlabel{Ib6}Elena trasformarsi in una secchia.
\pend


\pstart
Line 1 \\
Line 2 \\
Line 3 \\
Line 4 \\
Line 5 \\
Line 6 \\
Line 7 \\
\edtext{Line 8}{\Afootnote{A note on line 8, but not the same pstart}}\pend
\endnumbering

\noindent IB3 pstart : \xpstartref{Ib3}\\
IB4 pstart : \xpstartref{Ib4}

\makeatletter
\csuse{@this@pstart@end}
\csuse{@this@stanza@end}
\end{document}
